%% toc: P3 40 40.4
\chapter{Constant-Probability Regimes and Their Countable Models}
\label{sec:p03-w03-p03-067-constant-p}

This chapter treats the random structures whose basic probability does not
depend on the number of points: the random graph \(G(n,p)\) and the random
unary predicate \(SU(n,p)\) with \(0<p<1\) fixed.  For each it gives an
explicit first-order theory, proves that every axiom of the theory holds with
probability tending to one, proves the theory complete, and derives the
zero--one law.  It then turns from limits to models.  Whenever a zero--one law
holds, the sentences of limiting probability one form a complete theory with
infinite models only; in the constant-probability regimes that theory has
exactly one countable model up to isomorphism, and for graphs that model is
the Rado graph.  A final section shows where the constant-probability argument
breaks: once \(p(n)\) decays polynomially, some extension axiom holds almost
never, and in the sparse regime \(p=n^{-\alpha}\) with \(\alpha\) irrational the
almost sure theory, though still complete, has more than one countable model.

The inputs are a fixed probability \(p\) and a first-order sentence \(B\) in
the language of graphs (or of one unary predicate).  The conclusions are the
existence of \(\lim_{n\to\infty}\Pr[G(n,p)\entails B]\), its value \(0\) or
\(1\), its independence of \(p\), and a description of the countable models
of the almost sure theory, which determine the value.

The presentation follows Spencer's author version of ``Logic and Random
Structures'' \cite{Spencer2007LogicRandomStructuresAuthor}, a separately
formatted 30-page text whose sections are numbered 1 to 14 without a chapter
prefix; its page numbers are cited throughout.  It does not follow the
printed Springer chapter.  The constant-\(p\) graph law itself is due to
Glebski\u\i, Kogan, Liogon'ki\u\i\ and Talanov and, independently, to Fagin
\cite[p.~12]{Spencer2007LogicRandomStructuresAuthor}; a treatment along
Fagin's route is \cref{thm:p09-fagin1976-graph}.

\section{Random structures and the almost sure theory}
\label{sec:p03-w03-p03-067-setting}

\begin{hbconvention}{def:p03-w03-p03-067-assumptions}{Standing assumptions}
Every language here is finite and relational, with equality: either one
unary predicate \(U\), or one binary adjacency symbol \(\sim\).  A random
structure of size \(n\) is a probability distribution on structures whose
universe is a fixed set of exactly \(n\) points, for every sufficiently large
positive integer \(n\).  A first-order sentence \(A\) holds \emph{almost
surely} if \(\lim_{n\to\infty}\Pr[A]=1\) and \emph{almost never} if this limit
is \(0\), equivalently if \(\neg A\) holds almost surely
\cite[p.~2]{Spencer2007LogicRandomStructuresAuthor}.  The probability \(p\) is
a constant with \(0<p<1\) unless a result says that \(p=p(n)\) varies with
\(n\).  The metatheory is ZFC; choice enters only through compactness and the
downward L\"owenheim--Skolem theorem.
\end{hbconvention}

\begin{hbdefinition}{def:p03-w03-p03-067-models}{The two constant-probability models}
Let \(n\geq1\) and \(0\leq p\leq1\).
\begin{itemize}
\item \(SU(n,p)\), the \emph{simple unary predicate}, has a universe
  \(\Omega\) of \(n\) points; each \(x\in\Omega\) satisfies \(U(x)\) with
  probability \(p\), and these \(n\) events are mutually independent
  \cite[p.~1]{Spencer2007LogicRandomStructuresAuthor}.  The universe carries
  no further structure; in particular it is not ordered.
\item \(G(n,p)\), the \emph{random graph}, has a vertex set \(V\) of \(n\)
  points; adjacency \(\sim\) is irreflexive and symmetric; for each pair of
  distinct \(x,y\), \(\Pr[x\sim y]=p\), and these \(\binom n2\) events are
  mutually independent \cite[p.~4]{Spencer2007LogicRandomStructuresAuthor}.
\end{itemize}
\end{hbdefinition}

The regime is \emph{constant} when \(p\) does not depend on \(n\).  Spencer
lets \(p=p(n)\) vary in general and locates the thresholds at which a property
changes from almost never to almost surely
\cite[pp.~2, 4]{Spencer2007LogicRandomStructuresAuthor}.  Each constant
\(p\in(0,1)\) sits above every threshold function that tends to zero, which is
why one theory serves for all of them.

\begin{hbtheorem}{thm:p03-w03-p03-067-as-theory}{proved=annex:p03-w03-p03-067-as-theory}{The almost sure theory of a zero--one law}
Fix a finite relational language with equality and a random structure of
size \(n\) in the sense of \cref{def:p03-w03-p03-067-assumptions}.  Suppose
every first-order sentence holds almost surely or almost never, and let
\(T_{\mathrm{as}}\) be the set of sentences that hold almost surely.  Then
\begin{enumerate}
\item \(T_{\mathrm{as}}\) is consistent, closed under logical consequence and
  complete;
\item every model of \(T_{\mathrm{as}}\) is infinite;
\item \(T_{\mathrm{as}}\) has a countably infinite model.
\end{enumerate}
\end{hbtheorem}

This is the opening observation of Spencer's Section~7
\cite[p.~13]{Spencer2007LogicRandomStructuresAuthor}.  It turns the study of
limiting probabilities into the study of the countable models of one complete
theory.  It says nothing about which theory arises; for that one needs an
explicit axiomatization, given next for graphs.

\section{The extension schema for constant probability}
\label{sec:p03-w03-p03-067-extension}

\begin{hbdefinition}{def:p03-w03-p03-067-schema}{Extension axioms and the theory \(T\)}
For integers \(r,s\geq0\), the \emph{extension axiom} \(A_{r,s}\) is the
graph sentence
\begin{align*}
 &\forall x_1\cdots\forall x_r\,\forall y_1\cdots\forall y_s
 \Bigl[\mathrm{Dist}(\bar x,\bar y)\to\\
 &\qquad\exists z\Bigl(\bigwedge_{i\leq r}(z\neq x_i\land z\sim x_i)
 \land\bigwedge_{j\leq s}(z\neq y_j\land\neg\, z\sim y_j)\Bigr)\Bigr],
\end{align*}
where \(\mathrm{Dist}(\bar x,\bar y)\) says that the \(r+s\) variables take
pairwise distinct values.  In words: for all distinct
\(x_1,\ldots,x_r,y_1,\ldots,y_s\) there is a distinct \(z\) adjacent to every
\(x_i\) and to no \(y_j\).  Given the \(x_i\) and \(y_j\), a vertex \(z\) with
exactly these adjacencies is a \emph{witness}.

\(T\) is the theory whose axioms are irreflexivity and symmetry of \(\sim\)
together with \(A_{r,s}\) for all \(r,s\geq0\).  A graph satisfying every
\(A_{r,s}\) is said to have the \emph{Alice's Restaurant property}, a name due
to Winkler: every possible witness is present
\cite[pp.~12--13]{Spencer2007LogicRandomStructuresAuthor}.
\end{hbdefinition}

Spencer states the schema alone; the graph axioms are added here because
models of \(T\) are compared as graphs.  They hold in every outcome of
\(G(n,p)\), so they do not affect any limit.  The axiom \(A_{0,0}\) says only
that the universe is nonempty.

\begin{hbtheorem}{thm:p03-w03-p03-067-witness}{proved=annex:p03-w03-p03-067-witness}{Every extension axiom holds almost surely}
Let \(0<p<1\) be constant, \(r,s\geq0\), \(\varepsilon:=p^r(1-p)^s\) and
\(n\geq r+s\).  Then
\[
 \Pr\bigl[G(n,p)\not\entails A_{r,s}\bigr]\;\leq\;
 n^{r+s}(1-\varepsilon)^{\,n-r-s},
\]
and the right-hand side tends to \(0\) as \(n\to\infty\).  Hence every
sentence of \(T\) holds almost surely in \(G(n,p)\).
\end{hbtheorem}

This is Spencer's estimate \cite[p.~13]{Spencer2007LogicRandomStructuresAuthor}.
For fixed roots the witness events of different candidates involve disjoint
sets of pairs, hence are independent; there are at most \(n^{r+s}\) choices of
roots; and for fixed \(r,s,p\) the exponential factor beats the polynomial one.
The constant \(\varepsilon\) is where constancy of \(p\) is used: it must stay
bounded away from zero.

\begin{hbconvention}{def:p03-w03-p03-067-game}{The game \(\mathrm{EHR}(G,G';t)\)}
Spencer writes \(\mathrm{EHR}(G,G';t)\) for the \(t\)-round
Ehrenfeucht game on graphs \(G,G'\)
\cite[p.~26]{Spencer2007LogicRandomStructuresAuthor}.  It is the game of
\cref{def:p03-ef-game} for the language \(\{\sim\}\): in each round Spoiler
picks a vertex of either graph and Duplicator answers in the other, and
Duplicator wins when the pairing of chosen vertices preserves equality,
adjacency and nonadjacency.  By \cref{thm:p03-ef}, Duplicator wins
\(\mathrm{EHR}(G,G';t)\) exactly when \(G\) and \(G'\) satisfy the same
sentences of quantifier rank at most \(t\).
\end{hbconvention}

\begin{hbtheorem}{thm:p03-w03-p03-067-complete}{proved=annex:p03-w03-p03-067-complete}{Duplicator stays alive; \(T\) is complete}
If \(G\) and \(G'\) are models of \(T\), of any cardinality, then Duplicator
wins \(\mathrm{EHR}(G,G';t)\) for every \(t\geq1\).  Consequently \(T\) is
consistent and complete.
\end{hbtheorem}

Duplicator's strategy is ``staying alive''
\cite[p.~13]{Spencer2007LogicRandomStructuresAuthor}: when Spoiler picks a
new vertex, Duplicator looks in the other graph for a new vertex with the same
adjacencies to the vertices already picked, and the Alice's Restaurant
property guarantees one.  No countable model and no isomorphism is needed for
completeness.

\begin{hbcorollary}{thm:p03-w03-p03-067-zero-one}{proved=annex:p03-w03-p03-067-zero-one}{Zero--one law for constant \(p\)}
Let \(0<p<1\) be constant.  For every first-order graph sentence \(B\), the
limit \(\lim_{n\to\infty}\Pr[G(n,p)\entails B]\) exists; it is \(1\) if
\(T\vdash B\) and \(0\) otherwise.  The limit does not depend on \(p\), and the
almost sure theory of \(G(n,p)\) is the deductive closure of \(T\).
\end{hbcorollary}

\begin{hbmethod}{def:p03-w03-p03-067-procedure}{Deciding a limit in the constant regime}
To find \(\lim_n\Pr[G(n,p)\entails B]\) for a first-order graph sentence
\(B\) and constant \(p\in(0,1)\):
\begin{enumerate}
\item Decide whether \(B\) holds in the countable model of \(T\)
  (\cref{thm:p03-w03-p03-067-categorical}); the limit is \(1\) if it does and
  \(0\) if not.  Equivalently, find a deduction of \(B\) or of \(\neg B\) from
  finitely many axioms of \(T\); by \cref{thm:p03-w03-p03-067-complete} one
  of the two exists.
\item For an explicit rate, list the axioms \(A_{r_1,s_1},\ldots,A_{r_m,s_m}\)
  used in that deduction.  Then
  \(\Pr[G(n,p)\text{ fails the conclusion}]\leq\sum_{k\leq m}
  n^{r_k+s_k}(1-p^{r_k}(1-p)^{s_k})^{n-r_k-s_k}\) for
  \(n\geq\max_k(r_k+s_k)\), by \cref{thm:p03-w03-p03-067-witness} and the
  union bound.
\end{enumerate}
Side conditions: \(B\) is first-order in the language \(\{\sim\}\) with
equality; \(p\) does not depend on \(n\) and lies strictly between \(0\) and
\(1\); graphs are labeled on exactly \(n\) vertices.  Step~2 depends on the
deduction chosen, and nothing here bounds the rate uniformly in \(B\).
\end{hbmethod}

\section{The countable models}
\label{sec:p03-w03-p03-067-countable}

\Cref{thm:p03-w03-p03-067-as-theory} gives the almost sure theory a countable
model.  Spencer's first question about it is whether it is unique, that is,
whether the theory is \(\aleph_0\)-categorical
\cite[p.~13]{Spencer2007LogicRandomStructuresAuthor}.

\begin{hbtheorem}{thm:p03-w03-p03-067-categorical}{proved=annex:p03-w03-p03-067-categorical}{The constant-\(p\) theory is \(\aleph_0\)-categorical}
\(T\) has no finite model: a graph with exactly \(m\geq1\) vertices fails
\(A_{m,0}\).  \(T\) has a countably infinite model, and any two countable
models of \(T\) are isomorphic.  The countable model is the graph of
\cref{thm:p03-random-graph}, the Rado graph, whose extension property is the
conjunction of all the \(A_{r,s}\).
\end{hbtheorem}

The isomorphism is built by Spencer's alternation of Left and Right Stages:
with both vertex sets labeled by the positive integers, a Left Stage gives the
least unmatched vertex of the first graph an image with the right adjacencies
to all matched vertices, a Right Stage does the same in the other direction,
and after stage \(2n\) the vertices labeled \(1,\ldots,n\) are matched on both
sides \cite[pp.~13--14]{Spencer2007LogicRandomStructuresAuthor}.  The name
``Rado graph'' is Spencer's \cite[p.~14]{Spencer2007LogicRandomStructuresAuthor};
Rado's own construction and its extension property are in
\cref{thm:p03-rado64-extension}.

\begin{hbcorollary}{thm:p03-w03-p03-067-rado-limit}{proved=annex:p03-w03-p03-067-rado-limit}{Limits are truth values in the Rado graph}
Let \(R\) be the Rado graph and \(0<p<1\) constant.  For every first-order
graph sentence \(B\),
\(\lim_{n\to\infty}\Pr[G(n,p)\entails B]=1\) if \(R\entails B\), and \(=0\)
if \(R\entails\neg B\).
\end{hbcorollary}

The unary constant regime has the same shape with a simpler model.

\begin{hbtheorem}{thm:p03-w03-p03-067-unary}{proved=annex:p03-w03-p03-067-unary}{The unary regime and its countable model}
Let \(T^{SU}\) consist of the sentences \(U_r\): there exist distinct
\(x_1,\ldots,x_r\) with \(U(x_i)\) for all \(i\), and \(N_r\): there exist
distinct \(x_1,\ldots,x_r\) with \(\neg U(x_i)\) for all \(i\), for every
\(r\geq1\).  Let \(p=p(n)\) satisfy \(np(n)\to\infty\) and
\(n(1-p(n))\to\infty\); every constant \(p\in(0,1)\) does.  Then
\begin{enumerate}
\item every sentence of \(T^{SU}\) holds almost surely in \(SU(n,p(n))\);
\item \(T^{SU}\) is consistent and complete, so every first-order sentence
  \(A\) in \(\{U\}\) has \(\lim_n\Pr[SU(n,p(n))\entails A]\in\{0,1\}\), with a
  value depending only on \(A\) and not on \(p(n)\) within the range;
\item \(T^{SU}\) is \(\aleph_0\)-categorical: its countable models are
  exactly the countable sets in which \(U\) and its complement are both
  infinite, and any two are isomorphic.
\end{enumerate}
\end{hbtheorem}

Spencer states parts 1 and 2 for this range of \(p(n)\) and calls the theory
\(T\) \cite[p.~3]{Spencer2007LogicRandomStructuresAuthor}; the superscript in
\(T^{SU}\) only separates it from the graph theory of
\cref{def:p03-w03-p03-067-schema}.  Part~3 is his remark that this theory is
\(\aleph_0\)-categorical \cite[p.~4]{Spencer2007LogicRandomStructuresAuthor}.
For parts 1 and 2 Spencer indicates only the binomial argument and the
Duplicator strategy; the annex proof supplies a Chebyshev estimate for part~1.

\begin{hbexample}{def:p03-w03-p03-067-k4}{A four-clique at constant and at sparse density}
Let \(B\) say that there are four distinct pairwise adjacent vertices.  From
\(A_{0,0}\) pick \(z_1\); from \(A_{1,0}\) pick \(z_2\neq z_1\) adjacent to
\(z_1\); from \(A_{2,0}\) pick \(z_3\) adjacent to both; from \(A_{3,0}\) pick
\(z_4\) adjacent to all three.  So \(A_{0,0},\ldots,A_{3,0}\vdash B\), and by
\cref{def:p03-w03-p03-067-procedure}, for \(p=1/2\) and \(n\geq3\),
\[
 \Pr[G(n,\tfrac12)\not\entails B]\leq
 n\bigl(\tfrac12\bigr)^{n-1}+n^2\bigl(\tfrac34\bigr)^{n-2}
 +n^3\bigl(\tfrac78\bigr)^{n-3},
\]
the \(A_{0,0}\) term being \(0\).  At \(n=100\) the bound exceeds \(1\) and says
nothing; at \(n=200\) it is about \(3.0\times10^{-5}\), almost all of it from the
\(A_{3,0}\) term.  The limit is \(1\) for every constant \(p\).

Now let \(p=n^{-\alpha}\) with \(\alpha>2/3\).  There are \(\binom n4\leq n^4/24\)
four-sets, each a clique with probability \(p^6\), so the expected number of
four-cliques is at most \(n^{4-6\alpha}/24\to0\), and by Markov's inequality
\(B\) holds almost never.  Spencer gives \(n^{-2/3}\) as the threshold
function of this property \cite[p.~4]{Spencer2007LogicRandomStructuresAuthor}.
The limit of one sentence thus depends on the regime even though, within
the constant regime, it does not depend on \(p\).
\end{hbexample}

\section{Where the constant-probability argument fails}
\label{sec:p03-w03-p03-067-limits}

The estimate of \cref{thm:p03-w03-p03-067-witness} needs
\(\varepsilon=p^r(1-p)^s\) bounded away from zero.  When \(p\) decays like a
power of \(n\), the conclusion itself fails.

\begin{hbtheorem}{thm:p03-w03-p03-067-sparse-failure}{proved=annex:p03-w03-p03-067-sparse-failure}{Polynomially small \(p\) defeats an extension axiom}
Let \(p=p(n)\) satisfy \(p(n)\leq n^{-\alpha}\) for all large \(n\), where
\(\alpha>0\) is fixed, and let \(r\) be an integer with \(r>1/\alpha\).  Then
\(A_{r,0}\) holds almost never in \(G(n,p(n))\).  In particular, for
\(p=n^{-\alpha}\) with \(\alpha>1/3\), the axiom \(A_{3,0}\), ``every three
vertices have a common neighbor'', holds almost never.
\end{hbtheorem}

For \(\alpha=\pi/7\approx0.448\), Spencer's running example, this is his
observation that three random vertices almost surely have no common neighbor
\cite[pp.~7--8]{Spencer2007LogicRandomStructuresAuthor}.  So \(T\) is not
the almost sure theory of any polynomially sparse regime, and the Rado graph
is not its countable model.  The sparse regime has its own almost sure
theory \(T_\alpha\), built from rooted graphs and closures.

\begin{hbconvention}{def:p03-w03-p03-067-sparse-notions}{Sparse-regime notions used in the comparison}
Fix an irrational \(\alpha\in(0,1)\).  Rooted graphs \((R,H)\), their
parameters \(r,v,e\), extensions over a tuple of vertices, and the notions
dense, sparse, rigid and safe are those of
\cref{def:p03-w03-p03-065-rooted-graph,def:p03-w03-p03-065-dense-sparse}.
For a set \(S\) of vertices of a graph \(G\) and a finite nonempty set \(N\)
disjoint from \(S\), \(e_G(N;S)\) is the number of edges of \(G\) with an end
in \(N\) and both ends in \(S\cup N\) (\cref{def:p03-w03-p03-068-standing});
rigidity of a vertex set over another and the \(t\)-closure
\(\mathrm{cl}_t(X)\) are those of
\cref{def:p03-w03-p03-068-rigid-over,def:p03-w03-p03-068-closure}.  The
\emph{closure} \(\mathrm{cl}(X)\) is the union of the \(\mathrm{cl}_t(X)\)
over all \(t\geq0\); it is not a first-order notion
\cite[pp.~14--15]{Spencer2007LogicRandomStructuresAuthor}.  \(T_\alpha\) is
the set of sentences holding almost surely in \(G(n,n^{-\alpha})\); its
axioms, Nonexistence and Generic Extension, are those of
\cref{def:p03-w03-p03-068-axioms} \cite[p.~10]{Spencer2007LogicRandomStructuresAuthor},
and its completeness is \cref{thm:p03-w03-p03-065-shelah-spencer} with
\cref{thm:p03-w03-p03-067-as-theory}.  Throughout this section
\(k=\lfloor1/\alpha\rfloor\) and \(c=1/\alpha-k\); since \(1/\alpha\) is
irrational and exceeds \(1\), \(k\geq1\) and \(c\in(0,1)\) is irrational.
\end{hbconvention}

That the two schemas determine \(T_\alpha\) is stated by Spencer and recorded
as \cref{thm:p03-w03-p03-068-axiomatized}.  The constructions below use it
in the following form, which reduces it to results proved in the handbook:
the almost sure validity of every axiom and the completeness of the axioms
(\cref{thm:p03-w03-p03-070-as-holds,thm:p03-w03-p03-070-complete}), which
are established there for the axioms written with weights in
\cref{def:p03-w03-p03-070-theory}.  What is added here is that the axioms of
\cref{def:p03-w03-p03-068-axioms} are the same axioms.

\begin{hblemma}{thm:p03-w03-p03-067-axioms-generate}{proved=annex:p03-w03-p03-067-axioms-generate}{The two schemas axiomatize \(T_\alpha\)}
Let \(\alpha\in(0,1)\) be irrational.  Let \(\Sigma_\alpha\) consist of the
graph axioms (\(\sim\) is irreflexive and symmetric) together with the
Nonexistence and Generic Extension axioms of
\cref{def:p03-w03-p03-068-axioms}, where a copy of a finite graph \(H\) in
\(G\) is an injective map \(V(H)\to V(G)\) carrying edges of \(H\) to edges
of \(G\).
\begin{enumerate}
\item For every graph \(G\) and every datum, either a finite graph \(H\) with
  \(v-e\alpha<0\) or a safe rooted graph \((R,H)\) with an integer
  \(t\geq0\), \(G\) satisfies the axiom of
  \cref{def:p03-w03-p03-068-axioms} for that datum if and only if it
  satisfies the axiom of \cref{def:p03-w03-p03-070-theory} for the same
  datum.  Hence \(\Sigma_\alpha\) and the theory \(T^-_\alpha\) of
  \cref{def:p03-w03-p03-070-theory} have the same models.
\item Every sentence of \(\Sigma_\alpha\) holds almost surely in
  \(G(n,n^{-\alpha})\).
\item A first-order sentence holds almost surely in \(G(n,n^{-\alpha})\) if
  and only if it is a logical consequence of \(\Sigma_\alpha\).  So
  \(T_\alpha\) is the set of logical consequences of \(\Sigma_\alpha\), it is
  consistent and complete, and a structure for \(\{\sim\}\) is a model of
  \(T_\alpha\) if and only if it is a graph satisfying every axiom of
  \cref{def:p03-w03-p03-068-axioms}.
\end{enumerate}
\end{hblemma}

Item~3 is the content of \cref{thm:p03-w03-p03-068-axiomatized}.  Its
probabilistic part, the large-deviation estimate behind Generic Extension
that the source does not prove
\cite[pp.~10--11]{Spencer2007LogicRandomStructuresAuthor}, enters only
through \cref{thm:p03-w03-p03-070-as-holds}, which is proved in the handbook
from \cref{thm:p03-w03-p03-064-theory-as}; completeness enters through the
Duplicator strategy of \cref{thm:p03-w03-p03-070-complete}.  An alternative
route to the same completeness is the closure strategy of
\cref{thm:p03-w03-p03-072-sparse-zero-one}.

Spencer gives two constructions of countable models of \(T_\alpha\) and
states their properties with the proofs omitted
\cite[pp.~14--15]{Spencer2007LogicRandomStructuresAuthor}.  The first needs
rigid one-root extensions whose parameter \(v-e\alpha\) is negative but
arbitrarily close to \(0\); Spencer says only that a little number theory
produces such a sequence for every \(\alpha\)
\cite[p.~14]{Spencer2007LogicRandomStructuresAuthor}.  The arithmetic lemma
and the circulant construction below are written here to supply that step;
they are not taken from the source.

\begin{hblemma}{thm:p03-w03-p03-067-diophantine}{proved=annex:p03-w03-p03-067-diophantine}{Ceilings just above an irrational multiple}
Let \(c\in(0,1)\) be irrational, \(\eta>0\) and \(M\geq1\) an integer.  There
is an integer \(v\geq M\) with \(0<\lceil vc\rceil-vc<\eta\).
\end{hblemma}

\begin{hbdefinition}{def:p03-w03-p03-067-circulant}{The rooted graphs \((\{x\},H(v))\)}
Let \(v\geq2k+1\).  The graph \(H(v)\) has the root \(x\) and the nonroots
\(1,\ldots,v\).  Nonroots \(i\neq j\) are adjacent when their cyclic distance
\(\min(\abs{i-j},v-\abs{i-j})\) is at most \(k\), so the nonroots carry a
\(2k\)-regular circulant graph \(K(v)\) with \(kv\) edges.  The root is
adjacent exactly to the nonroots in
\[
 Q_v=\{\,j\in\{1,\ldots,v\}:\lceil jc\rceil>\lceil (j-1)c\rceil\,\},
\]
a set of \(\lceil vc\rceil\) elements.  So \((\{x\},H(v))\) has parameters
\(v\) and \(e=kv+\lceil vc\rceil\).
\end{hbdefinition}

\begin{hbtheorem}{thm:p03-w03-p03-067-circulant}{proved=annex:p03-w03-p03-067-circulant}{Rigid one-root extensions of nearly zero excess}
Let \(v\geq2k+1\) and \(g(v)=\lceil vc\rceil-vc\).
\begin{enumerate}
\item \(v-e\alpha=-\alpha\,g(v)<0\).
\item \((\{x\},H(v))\) is rigid.
\item Every nonempty set \(S\) of nonroots satisfies
  \(\abs S-e(K(v)|_S)\,\alpha>0\).
\item Let \(v_1,v_2,\ldots\geq2k+1\) satisfy \(g(v_i)<2^{-i}\), which
  \cref{thm:p03-w03-p03-067-diophantine} makes possible, and let \(H^s\)
  be the union of \(H(v_1),\ldots,H(v_s)\) with disjoint nonroots and the
  common root \(x\).  Then \((\varnothing,H^s)\) is safe for every
  \(s\geq1\), and almost surely \(G(n,n^{-\alpha})\) contains a copy of
  \(H^s\).
\end{enumerate}
\end{hbtheorem}

These are the \(H_i=H(v_i)\) and the glued graphs \(H^s\) of Spencer's first
construction \cite[p.~14]{Spencer2007LogicRandomStructuresAuthor}.  The
rooted graph \((\{x\},H(v))\) is one rigid step of size \(v\); its excess
\(\alpha g(v)\) is what a single rigid step costs, and the condition
\(g(v_i)<2^{-i}\) keeps the total cost of all steps below \(1\), which is the
budget of the one vertex \(x\).

\begin{hbtheorem}{thm:p03-w03-p03-067-infinite-closure}{proved=annex:p03-w03-p03-067-infinite-closure}{A countable model with an infinite closure}
\(T_\alpha\) has a countable model containing a vertex \(x\) with
\(\mathrm{cl}(\{x\})\) infinite.
\end{hbtheorem}

The model comes from compactness: a new constant is required to be the root
of a copy of \(H^s\) for every \(s\)
\cite[p.~14]{Spencer2007LogicRandomStructuresAuthor}.  This does not
conflict with the Finite Closure Lemma
(\cref{thm:p03-w03-p03-068-finite-closure}), which bounds each
\(t\)-closure separately: the argument places a copy of \(H(v_i)\) at \(x\)
inside \(\mathrm{cl}_t(\{x\})\) only for \(t\geq v_i\), and the \(v_i\)
cannot stay bounded, because each fixed \(v\) has a fixed positive
\(g(v)\).

\begin{hbdefinition}{def:p03-w03-p03-067-demands}{The witness-demand construction}
Fix a set of representatives of the safe rooted graphs \((R,H)\) up to
isomorphism of rooted graphs with listed roots \(x_1,\ldots,x_r\) and
nonroots \(y_1,\ldots,y_v\); there are countably many, and \(r=0\) is
allowed.  A \emph{witness demand} is a pair of a representative \((R,H)\)
and a tuple \(\bar a=(a_1,\ldots,a_r)\) of distinct positive integers
\cite[p.~14]{Spencer2007LogicRandomStructuresAuthor}.  Fix an enumeration
of all witness demands.  A \emph{run} builds finite graphs
\(G_0\subseteq G_1\subseteq\cdots\) on vertex sets
\(\varnothing=P_0\subseteq P_1\subseteq\cdots\) of positive integers.  At
stage \(m\geq1\), take the first demand in the enumeration that has not been
met and whose roots \(a_1,\ldots,a_r\) all lie in \(P_{m-1}\); choose \(v\)
distinct positive integers \(b_1,\ldots,b_v\) outside \(P_{m-1}\), the
\emph{fresh set} \(F_m\); put \(P_m=P_{m-1}\cup F_m\); and add to
\(G_{m-1}\) the edges \(b_ib_j\) for \(y_iy_j\in E(H)\) and \(a_ib_j\) for
\(x_iy_j\in E(H)\), and no others.  The demand is then \emph{met} at stage
\(m\).  The graph \(G^\ast\) of the run is the union of the \(G_m\).
\end{hbdefinition}

Spencer turns the demands into one list and meets them in order with
untouched points, joining each new point only where the demand requires
\cite[p.~14]{Spencer2007LogicRandomStructuresAuthor}.  He does not say what
happens to a demand whose roots have not yet been touched; the rule above
postpones such a demand until its roots are present.  If each fresh set is
taken to be the least integers outside \(P_{m-1}\), then \(P_m\) contains
\(1,\ldots,m\), the vertex set of \(G^\ast\) is the set of all positive
integers, and every demand is eventually met, as in Spencer's description.

\begin{hbtheorem}{thm:p03-w03-p03-067-demand-model}{proved=annex:p03-w03-p03-067-demand-model}{The witness-demand graph models \(T_\alpha\) and has finite closures}
Let \(G^\ast\) be the graph of a run, with stage sets \(P_m\).
\begin{enumerate}
\item For every \(m\geq0\) and every finite nonempty set \(Y\) of vertices
  of \(G^\ast\) outside \(P_m\), \(\abs Y-e_{G^\ast}(Y;P_m)\,\alpha>0\).
\item Every demand whose roots are vertices of \(G^\ast\) is met, and
  \(G^\ast\) is countably infinite.
\item If a demand \(((R,H),\bar a)\) is met at stage \(m\) with fresh
  vertices \(b_1,\ldots,b_v\), then in \(G^\ast\) the \(b_j\) have no
  adjacencies among themselves or to the \(a_i\) other than those required
  by \(H\), and they are \(t\)-generic over \(a_1,\ldots,a_r\) for every
  \(t\geq1\).
\item \(G^\ast\) satisfies every axiom of \(T_\alpha\), and is a model of
  \(T_\alpha\).
\item Every finite set \(X\) of vertices of \(G^\ast\) has
  \(\mathrm{cl}(X)\subseteq P_m\) for any \(m\) with \(X\subseteq P_m\); in
  particular \(\mathrm{cl}(X)\) is finite.
\end{enumerate}
\end{hbtheorem}

Item~3 is Spencer's parenthetical expectation that minimal extensions are
\(t\)-generic for all \(t\), and items 4 and 5 are his two claims about the
construction \cite[pp.~14--15]{Spencer2007LogicRandomStructuresAuthor}.
The annex proves that \(G^\ast\) satisfies every axiom of
\cref{def:p03-w03-p03-068-axioms}; that a graph satisfying these axioms is a
model of \(T_\alpha\) is item~3 of
\cref{thm:p03-w03-p03-067-axioms-generate}.  The source proves
Nonexistence, sketches Generic Extension from a large-deviation estimate it
does not prove, and sketches completeness
\cite[pp.~10--12]{Spencer2007LogicRandomStructuresAuthor}; in the handbook
these steps are proved in the results that lemma rests on.

\begin{hbtheorem}{thm:p03-w03-p03-067-demand-unique}{proved=annex:p03-w03-p03-067-demand-unique}{The witness-demand graph is unique}
Any two runs, with any enumerations of the demands and any choices of fresh
sets, produce isomorphic graphs.
\end{hbtheorem}

Spencer finds this surprising and gives no proof
\cite[p.~14]{Spencer2007LogicRandomStructuresAuthor}.  The annex proof is a
back-and-forth between finite vertex sets over which every finite outside
set is sparse; the stage sets of a run are such sets by item~1 of
\cref{thm:p03-w03-p03-067-demand-model}.

\begin{hbcorollary}{thm:p03-w03-p03-067-sparse-models}{proved=annex:p03-w03-p03-067-sparse-models}{The sparse theory is not \(\aleph_0\)-categorical}
Let \(\alpha\in(0,1)\) be irrational.  \(T_\alpha\) has a countable model
containing a vertex \(x\) with \(\mathrm{cl}(\{x\})\) infinite, and a
countable model in which every finite vertex set has finite closure.  These
models are not isomorphic, so \(T_\alpha\) is not \(\aleph_0\)-categorical.
\end{hbcorollary}

This is Spencer's conclusion \cite[p.~15]{Spencer2007LogicRandomStructuresAuthor}.
The almost-sure-model classification therefore splits by regime.  In the
constant regimes the almost sure theory has exactly one countable model up
to isomorphism: the Rado graph for \(G(n,p)\)
(\cref{thm:p03-w03-p03-067-categorical}) and the set with two infinite
classes for \(SU(n,p)\) (\cref{thm:p03-w03-p03-067-unary}).  In the regime
\(p=n^{-\alpha}\) with \(\alpha\) irrational it has at least two,
distinguished by whether some vertex has infinite closure, and the
witness-demand graph is one of them, determined up to isomorphism by its
construction.  The source gives no complete list of the countable models of
\(T_\alpha\), and none is claimed here.

\begin{hbexample}{def:p03-w03-p03-067-pi-seven}{One rigid step at \(\alpha=\pi/7\)}
Spencer's running value is \(\alpha=\pi/7=0.44879\ldots\), and his example of
a rigid one-root extension \((\{x\},H)\) with \((\varnothing,H)\) safe is
\(H=K_5\) \cite[p.~14]{Spencer2007LogicRandomStructuresAuthor}: it has
\(v=4\), \(e=10\) and \(v-e\alpha=-0.488\ldots\), while every set of \(j\)
vertices of \(K_5\) has \(j-\binom j2\alpha>0\).  Gluing \(s\) copies of
\(K_5\) at \(x\) gives a graph whose parameter over \(\varnothing\) is
\(1-0.488\ldots s\), negative from \(s=3\) on; so copies of \(K_5\) cannot
serve as all the \(H_i\).

For the construction of \cref{def:p03-w03-p03-067-circulant},
\(1/\alpha=2.2281\ldots\), \(k=2\) and \(c=0.2281\ldots\).  At \(v=13\),
\(13c=2.9661\ldots\), so \(\lceil13c\rceil=3\) and
\(g(13)=0.0338\ldots\).  The nonroots \(1,\ldots,13\) form the square of a
\(13\)-cycle, each joined to the four nonroots at cyclic distance \(1\) or
\(2\); the root is joined to \(Q_{13}=\{1,5,9\}\), the indices at which
\(\lceil jc\rceil\) first takes the values \(1,2,3\).  Then \(e=26+3=29\) and
\(v-e\alpha=13-29\pi/7=-0.0151\ldots=-\alpha g(13)\).  One rigid step of
this kind costs less than one thirtieth of what a step of \(K_5\) costs, and
since \(g(13)<2^{-4}\), \(H(13)\) can serve as \(H_i\) for every
\(i\leq4\).
\end{hbexample}

\begin{hbwarning}{def:p03-w03-p03-067-scope-warning}{What the limit law does not say}
The zero--one law is asymptotic.  Every finite graph fails some axiom of
\(T\) (\cref{thm:p03-w03-p03-067-categorical}), so no finite graph is a model
of the almost sure theory, and limit one for \(B\) does not make \(B\) true in
\(G(n,p)\) for any particular \(n\).  The Rado graph is a model of the limiting
theory, not a limit of the finite graphs.  The law concerns first-order
sentences over \(\{\sim\}\) alone.  The independence of the limit from \(p\)
holds only among constants in \((0,1)\): by
\cref{thm:p03-w03-p03-067-sparse-failure} and
\cref{def:p03-w03-p03-067-k4}, the same sentence can have limit \(1\) at
every constant \(p\) and limit \(0\) at \(p=n^{-\alpha}\).
\end{hbwarning}

\section{Relationships and sources}
\label{sec:p03-w03-p03-067-relationships}

\Cref{sec:p03-zero-one-law} proves the extension estimate for \(G(n,1/2)\)
and cites the first-order zero--one law for uniformly random relational
structures; the present chapter replaces \(1/2\) by an arbitrary constant
\(p\) and proves completeness by the game.  \Cref{thm:p09-fagin1976-graph}
establishes the same constant-\(p\) graph law by a different route,
back-and-forth followed by an elementary-substructure argument; Spencer's
route needs no countable model for completeness and uses
\(\aleph_0\)-categoricity only to identify the model.  The Rado graph is
characterized deterministically in \cref{thm:p03-random-graph} and
constructed in \cref{thm:p03-rado64-extension}; the Ehrenfeucht game and its
rank theorem are \cref{def:p03-ef-game} and \cref{thm:p03-ef}.  For the
sparse regime, rooted graphs, extension statements and the safe criterion
are in \cref{sec:p03-w03-p03-065-extension-statements}; rigid extensions,
closures, the Finite Closure Lemma and the axioms of \(T_\alpha\) are in
\cref{sec:p03-w03-p03-068-closure}; this chapter uses them only to compare
the countable models of \(T_\alpha\) with those of the constant regimes.

Spencer's text supplies the definitions of \(SU(n,p)\) and \(G(n,p)\)
\cite[pp.~1, 4]{Spencer2007LogicRandomStructuresAuthor}, the unary theorem
and its countable models \cite[pp.~3--4]{Spencer2007LogicRandomStructuresAuthor},
the constant-\(p\) schema, witness estimate and Duplicator strategy
\cite[pp.~12--13]{Spencer2007LogicRandomStructuresAuthor}, the countable-model
analysis \cite[pp.~13--15]{Spencer2007LogicRandomStructuresAuthor}, and the
game \cite[pp.~26--27]{Spencer2007LogicRandomStructuresAuthor}.  Spencer
gives the almost-sure-theory statement as a remark and the unary estimate as
an appeal to the binomial distribution; the proofs of both, the sparse
failure theorem and the explicit four-clique bound are written here from
those ingredients and are not taken from his text.  For the sparse countable
models Spencer states the two constructions and their properties without
proof \cite[pp.~14--15]{Spencer2007LogicRandomStructuresAuthor}; the
arithmetic lemma, the circulant rooted graphs \(H(v)\), the rule for
postponing demands, and the proofs that the witness-demand graph satisfies
the axioms, has finite closures and is unique are written here.  The
axiomatization of \(T_\alpha\), stated by the source
(\cref{thm:p03-w03-p03-068-axiomatized}) with its large-deviation estimate
omitted \cite[pp.~10--11]{Spencer2007LogicRandomStructuresAuthor}, is used
in the form of \cref{thm:p03-w03-p03-067-axioms-generate}, which identifies
the axioms of \cref{def:p03-w03-p03-068-axioms} with those of
\cref{def:p03-w03-p03-070-theory} and draws on the proofs of almost sure
validity and completeness in
\cref{thm:p03-w03-p03-070-as-holds,thm:p03-w03-p03-070-complete}.
